\labsection{2}{Audit missingness, categories, and validity}{sec:lab02}

\begin{labproject}{Data-quality audit: define rules before cleaning}
\textbf{Coverage.} Chapter~\ref{ch:quality}.\\
\textbf{Goal.} You should be able to distinguish missingness, invalid values, inconsistent categories, duplicate records, and structural absence; then choose a defensible treatment based on variable meaning and analytical purpose.\\
\textbf{Datasets.} \texttt{iris\_dataset.csv} and \texttt{titanic\_dataset.csv}; Ames Housing and Red Wine Quality are used as extensions.\\
\textbf{Scenario.} You have received data from multiple sources and must decide whether it is fit for analysis. The objective is not to make the data ``look clean'', but to identify defects, document rules, and preserve an auditable record of every transformation.

\courselabflow{Profile fields}{Flag problems}{Document repairs}{Recheck}


\textbf{Run the complete script.} Start a clean R session in the course folder. Each step continues the same file; comment headings identify the operation performed by the next block.
\begin{lstlisting}[style=utpterminal]
Rscript all_experiments/lab02_data_quality.R
\end{lstlisting}

\noindent\textbf{Step 1.} Load, validate and prepare the data
\lstinputlisting[style=utpr,firstline=1,lastline=26]{all_experiments/lab02_data_quality.R}

\noindent\textbf{Step 2.} Standardise category labels without modifying the raw field.
\courselabtrace{Keep raw values recoverable}{Raw}{Retain the supplied columns unchanged}{Derived}{Write standardised or imputed values to new fields}{Audit}{Count changes and explain the rule used}

\lstinputlisting[style=utpr,firstline=27,lastline=48]{all_experiments/lab02_data_quality.R}

\noindent\textbf{Step 3.} Demonstrate a documented derived field rather than overwriting Age.


\textbf{Use the chapter while coding.} Revisit \hyperref[ch:quality]{the earlier chapter explanation}; match its operation to the variables and state changes below.

\textbf{Follow the supplied implementation.} Read the assignment and the condition around it before interpreting the count. This is a documented cleaning demonstration. If using imputation in a prediction workflow, learn its replacement value from training rows only.
\courselabtrace{Trace the derived age field}{1. Input / state}{Keep Titanic Age as the raw field and initialise imputed\_n for the audit.}{2. Operation}{Age\_imputed fills missing values using medians within Pclass and Sex groups; imputed\_n counts replacements.}{3. Check}{Compare raw missingness, imputed\_n and the derived values; raw Age must remain recoverable.}


\lstinputlisting[style=utpr,firstline=49,lastline=68]{all_experiments/lab02_data_quality.R}

\noindent\textbf{Step 4.} Extension A - Ames Housing: inspect high-missingness fields before assigning meaning.
\lstinputlisting[style=utpr,firstline=69,lastline=78]{all_experiments/lab02_data_quality.R}

\noindent\textbf{Step 5.} Extension B - Red Wine Quality: identical feature rows are not automatically bad records.
\lstinputlisting[style=utpr,firstline=79,lastline=85]{all_experiments/lab02_data_quality.R}
\textbf{Chapter connection.} Revisit \hyperref[ch:quality]{the chapter explanation} and map its states to the steps below.
% Reused from chapters/02-data-quality.tex
\begin{labvisual}{Recall the chapter model}
\begin{labgraphic}
\begin{tikzpicture}[x=\linewidth,y=1cm,
  sourcecard/.style={draw=UTPBlue!75,fill=UTPPaleBlue,rounded corners=1.8mm,
    line width=0.75pt,text width=0.15\linewidth,
    minimum width=0.19\linewidth,minimum height=1.15cm,
    inner xsep=3mm,inner ysep=2.5mm,align=center,
    font=\scriptsize\color{UTPNavy}},
  taskcard/.style={draw=UTPBlue!65,fill=UTPPaleBlue,rounded corners=1.8mm,
    line width=0.7pt,text width=0.21\linewidth,
    minimum width=0.255\linewidth,minimum height=1.05cm,
    inner xsep=3.5mm,inner ysep=2mm,align=left,
    font=\scriptsize\color{UTPNavy}},
  qualityflow/.style={-{Stealth[length=2mm,width=1.5mm]},
    line width=0.85pt,draw=UTPTeal}]
\node[sourcecard] (data) at (0.12,0)
  {\textcolor{UTPBlue}{\textbf{SHARED INPUT}}\\[3pt]
   The same dataset};
\node[taskcard] (q1) at (0.43,0.80)
  {\textcolor{UTPBlue}{\textbf{QUESTION A}}\\[3pt]
   Requires a complete target.};
\node[taskcard] (q2) at (0.43,-0.80)
  {\textcolor{UTPBlue}{\textbf{QUESTION B}}\\[3pt]
   Requires reliable dates.};
\node[taskcard,draw=UTPTeal!80,fill=UTPPaleTeal] (a1) at (0.78,0.80)
  {\textcolor{UTPTeal}{\textbf{QUALITY CHECK}}\\[3pt]
   Is the target complete and reliable?};
\node[taskcard,draw=UTPTeal!80,fill=UTPPaleTeal] (a2) at (0.78,-0.80)
  {\textcolor{UTPTeal}{\textbf{QUALITY CHECK}}\\[3pt]
   Are the dates complete and valid?};
\coordinate (split) at (0.275,0);
\draw[draw=UTPTeal,line width=0.85pt] (data.east) -- (split);
\draw[qualityflow] (split) |- (q1.west);
\draw[qualityflow] (split) |- (q2.west);
\draw[qualityflow] (q1.east) -- (a1.west);
\draw[qualityflow] (q2.east) -- (a2.west);
\fill[UTPTeal] (split) circle (1.3pt);
\end{tikzpicture}
\end{labgraphic}\end{labvisual}


\begin{verification}
\textbf{Explain and extend your result.} Audit one cleaning decision with before/after values and a reason. Compare the relevant distribution before and after treatment, including the count of observations.


\textbf{Required evidence.} Submit a data-quality log with each issue written as \emph{rule $\rightarrow$ evidence $\rightarrow$ action $\rightarrow$ consequence}. Include at least one issue you deliberately leave unchanged because the evidence is insufficient to justify a transformation.

\textbf{Professional standard.} A cleaning step is acceptable only when another analyst can determine what changed, why it changed, and whether the original value is still recoverable.
\end{verification}
\end{labproject}
